\subsection{Comparaciones conjuntas}

Se agrupa \var{tenure} en 0--12, 13--24, 25--48 y 49--72 meses para distinguir hasta un año, el 
segundo año, entre dos y cuatro años y más de cuatro años de permanencia. Estos intervalos se eligieron 
para facilitar la comparación y no fueron definidos por un modelo. La variable original se conserva sin 
modificar. Las proporciones se calculan dentro de cada combinación, aunque los intervalos tengan distinta 
amplitud.

La Fig.~\ref{fig:contrato_tiempo} muestra que los contratos mensuales presentan más abandono dentro de 
cada intervalo. En 0--12 meses, las tasas son 51.35\%, 10.48\% y 0\% para contratos mensuales, anuales y 
de dos años. Sin embargo, los tamaños respectivos son 1994, 124 y 68 clientes.  En el grupo de 68 clientes 
no se observaron abandonos, pero eso no significa que su probabilidad real sea cero. En 49--72 meses las 
tasas son 26.02\%, 12.93\% y 3.33\%, con 342, 634 y 1263 clientes.

\begin{figure}[t]
    \centering\includegraphics[width=\columnwidth]{contrato_antiguedad.png}
    \caption{Porcentaje de abandono dentro de cada combinación de contrato y antigüedad. 
    Cada intervalo corresponde a un grupo distinto de clientes.}
    \label{fig:contrato_tiempo}
\end{figure}

La diferencia entre contratos se mantiene dentro de estos grupos de antigüedad. Aun así, los intervalos 
son amplios y quedan otras características sin comparar. Tampoco puede interpretarse que permanecer más 
meses reduzca por sí mismo el riesgo de abandono.

Al separar por contrato, la fibra también mantiene mayor abandono que DSL (Fig.~\ref{fig:contrato_internet}). 
Las tasas son 54.61\% frente a 32.22\% en contratos mensuales, 19.29\% frente a 9.30\% en anuales y 
7.23\% frente a 1.91\% en contratos de dos años. Las combinaciones de fibra y DSL cuentan, respectivamente, 
con 2128 y 1223 clientes mensuales, 539 y 570 anuales, y 429 y 628 de dos años. La diferencia no desaparece 
al separar por contrato, aunque todavía podrían intervenir cargos, antigüedad u otros atributos.

\begin{figure}[t]
    \centering\includegraphics[width=\columnwidth]{contrato_internet.png}
    \caption{Abandono según servicio de internet dentro de cada contrato. La categoría sin internet se 
    distingue de DSL y fibra.}
    \label{fig:contrato_internet}
\end{figure}
